



































































































































































































































































































































































































































































































































































\section{Introduction}
For a (labeled) graph $G=(V,E)$, and a bijection $\phi:V\to V'$,
we denote by $\phi(G)$ the graph $G'=(V',E')$
such that $E'=\{\{\phi(u),\phi(v)\}:\{u,v\}\!\in\!E\}$,
and say that $G'$ is {\sf isomorphic} to~$G$.
The set of {\sf automorphisms} of the graph $G=(V,E)$,
denoted ${\tt aut}(G)$,
is the set of permutations that preserve the graph~$G$;
that is, $\pi\in{\tt aut}(G)$ if and only if $\pi(G)=G$.
We say that a graph is {\sf asymmetric} (equiv., {\sf self-ordered})
if its set of automorphisms is a singleton,
which consists of the trivial automorphism
(i.e., the identity permutation).
We actually prefer the term {\em self-ordered}, because we take
the perspective that is offered by the following equivalent definition.

%
\BD[Self-ordered (a.k.a asymmetric) graphs]
\label{asymmetric:def}
%
The graph $G=([n],E)$ is {\sf self-ordered}
if for every graph $G'=(V',E')$ that is isomorphic to~$G$
there exists a unique bijection $\phi:V'\to[n]$ such that $\phi(G')=G$.
\ED
In other words,
given an isomorphic copy $G'=(V',E')$ of a fixed graph $G=([n],E)$,
there is a unique bijection $\phi:V'\to[n]$ that orders the vertices of $G'$
such that the resulting graph (i.e., $\phi(G')$) is identical to~$G$.
Indeed, if $G'=G$, then this unique bijection is the identity permutation.
\footnote{Naturally, we are interested in efficient algorithms
that find this unique ordering, whenever it exists;
such algorithms are known
when the degree of the graph is bounded~\cite{L}.}

In this work, we consider a feature,
which we call {\em robust self-ordering},
that is a quantitative version self-ordering.
%
Loosely speaking, a graph $G=([n],E)$ is robustly self-ordered if,
%
%
for every permutation $\pi:[n]\to[n]$,
the size of the symmetric difference between~$G$ and $\pi(G)$
is proportional to the number of non-fixed-points under $\pi$;
that is, $|E\triangle\{\{\pi(u),\pi(v)\}\!:\!\{u,v\}\!\in\!E\}|$
is proportional to $|\{i\!\in\![n]\!:\!\pi(i)\!\neq\!i\}|$.
(In contrast, self-ordering only means that
the size of the symmetric difference is positive
if the number of non-fixed-points is positive.)

%
\BD[Robustly self-ordered graphs]
\label{robust-asymmetric:def}
A graph $G=(V,E)$ is said to be {\sf $\gamma$-robustly self-ordered}
if for every permutation $\pi:V\to{V}$
it holds that
\begin{equation}\label{robust-asymmetric:eqdef}
\big|E\triangle\left\{\{\pi(u),\pi(v)\}\!:\!\{u,v\}\!\in\!E\right\}\big|
  \;\geq\; \gamma\cdot|\{i\in[n]\!:\!\pi(i)\neq i\}|,
\end{equation}
where $\triangle$ denotes the symmetric difference operation.
An infinite family of graphs $\{G_n=([n],E_n)\}_{n\in\N}$
{\em(such that each $G_n$ has maximum degree $d$)}
is called {\sf robustly self-ordered}
if there exists a constant $\gamma>0$, called the {\sf robustness parameter},
such that for every $n$ the graph $G_n$ is $\gamma$-robustly self-ordered.
\ED
Note that
$|E_n\triangle\{\{\pi(u),\pi(v)\}\!:\!\{u,v\}\!\in\!E_n\}|
 \leq 2d\cdot|\{i\in[n]\!:\!\pi(i)\neq i\}|$
always holds (for families of maximum degree $d$).
%
%
%
%
%
The term ``robust'' is inspired by
the property testing literature (cf.~\cite{RS}),
where it indicates that some ``parametrized violation''
is reflected proportionally in some ``detection parameter''.

The second part of Definition~\ref{robust-asymmetric:def}
is tailored for bounded-degree graphs, which will be our focus
in Section~2--\ref{random-graphs:sec}.
Nevertheless,
in Sections~7--10
we consider graphs
of unbounded degree and unbounded robustness parameters.
In this case, for a function $\rho:\N\to\R$,
we say that an infinite family of graphs $\{G_n=([n],E_n)\}_{n\in\N}$
is {\sf $\rho$-robustly self-ordered} if for every $n$
the graph $G_n$ is $\rho(n)$-robustly self-ordered.
Naturally, in this case, the graphs
must have $\Omega(\rho(n)\cdot n)$ edges.
\footnote{Actually, all but at most one vertex
must have degree at least $\rho(n)/2$.}
%
In Sections~7--9
we consider the case of $\rho(n)=\Omega(n)$.




































































































































































































































































































































































































































































































































































































































































































































































































































































































































































































































































































































































































































































































































































































































































































































































































































































































































































































































































































































































































































































































































































































































































































































































































































































































































































































































































































































































































































































































































































































































































































































































































































































































































































































































































































































































































































































































































































































































































































































































































































































































































































































































































































































































































































































































































































































































































































































































































































































































\setcounter{section}{5}\section{Random Regular Graphs are Robustly Self-Ordered}
%
\label{random-graphs:sec}
While Theorem~4.1 only asserts the
existence of robustly self-ordered $d$-regular graphs,
we next show that almost all $d$-regular graphs
are robustly self-ordered.
This extends work in probabilistic graph theory,
which proves a similar result for the weaker notion
of self-ordered (a.k.a asymmetric) graphs~\cite{B1,B2}.
%
%
%
%

\BT[Random $d$-regular graphs are robustly self-ordered]
\label{random-works:thm}
For any sufficiently large constant~$d$,
a random $2d$-regular $n$-vertex graph
is robustly self-ordered with probability $1-o(1)$.
%
%
\ET
Recall that, with very high probability, these graphs are expanders.
%
We mention that the proof of Theorem~4.1
actually established that $n$-vertex graphs drawn from
a weird distribution (which has min-entropy $\Omega(n)$)
are robustly self-ordered with probability $1-o(1)$.
However, this is established by using the edge-coloring variant,
and requires employing the transformation presented
in Section~2.1.
In contrast, the following proof works directly with the original
(uncolored) variant, and is completely self-contained.
\medskip
%
%
%

\BPF
The proof is quite similar to the proof Claim~4.1.1,
but it faces complications that were avoided in the prior proof
by using edge-colors and implicitly directed edges.
Specifically, for candidate permutations $\pi_1,\ldots,\pi_d:[n]\to[n]$
(to be used in the construction)
and all (non-trivial) permutations $\mu:[n]\to[n]$,
the proof of Claim~4.1.1 considered events
of the form $(\forall j\!\in\![d])\;\pi_j(i)=\mu(\pi_j(\imu(i)))$,
whereas here we shall consider events of the form
$\{\pi^b_j(i):j\!\in\![d]\,\&\,b\!\in\!\{\pm1\}\}
 =\{\mu(\pi_k^c(\imu(i))):k\!\in\![d]\,\&\,c\!\in\!\{\pm1\}\}$.
That is, rather than considering a sequence of equations that
refers to single $\pi_j$'s, we consider a sequence of equations
that refer to the set of all $\pi_j$'s, which were expressed
above in terms of equations between multi-sets.
These multi-set equations will be reduced to equations among sequences
by considering all possible ordering of these multi-sets.
This amounts to taking a union bound over all possible ordering
and results in a more complicated analysis (due to the $\pi_j^{-1}$'s)
and much more cumbersome notation.

To facilitate the proof,
we use the standard methodology (cf.~\cite[Apdx.~2]{E})
of first proving the result in the {\em random permutation model},
%
then transporting it to the {\em configuration model}\/
%
(by using a general result of~\cite{GJKW}),
and finally conditioning on the event that the generated graph
is simple (which occurs with positive constant probability).
Indeed, both models generate multi-graphs
that are not necessarily simple graphs
(i.e., these multi-graphs may have self-loops and parallel edges).
We also use the fact that the simple graphs that are
generated by the configuration model (for degree $d'$)
are uniformly distributed among all $d'$-regular graphs.

Recall that in the {\sf random permutation model}
a $2d$-regular $n$-vertex multi-graph is generated by selecting
uniformly and independently $d$ permutations $\pi_1,\ldots,\pi_d:[n]\to[n]$.
The multi-graph, denoted $G_{(\pi_1,\ldots,\pi_d)}$,
consists of the edge multi-set
$\bigcup_{j\in[d]}\{\{i,\pi_j(i)\}:i\in[n]\}$,
where the $2j^\xth$ (resp., $(2j-1)^\st$) neighbor of vertex $i$
is $\pi_j(i)$ (resp., $\pi_j^{-1}(i)$).
Note that this multi-graph may have self-loops (due to $\pi_j(i)=i$),
which contributed two units to the degree of a vertex,
as well as parallel edges
(due to $\pi_j(i)=\pi_k(i)$ for $j\neq k$
and $\pi_j(i)=\pi_k^{-1}(i)$ for any $j,k$).
We denote the $j^\xth$ neighbor of vertex $i$ by $g_j(i)$;
that is, $g_j(i)=\pi_{j/2}(i)$ if $j$ is even,
and $g_j(i)=\pi_{(j+1)/2}^{-1}(i)$ otherwise.
%

Consider an arbitrary permutation $\mu:[n]\to[n]$,
and let $T=\{i\!\in\![n]\!:\!\mu(i)\neq i\}$
be its set of non-fixed-point.
We shall show that,
with probability $1-\exp(-\Omega(d\cdot|T|\cdot\log n))$
over the choice of ${\ov\pi}=(\pi_1,\ldots,\pi_d)$,
the size of the symmetric difference
between $G_{\ov\pi}$ and $\mu(G_{\ov\pi})$ is $\Omega(|T|)$.
Note that this difference is (half) the sum over $i\in[n]$
of the size of the symmetric difference between the multi-set
of neighbors of vertex~$i$ in $G_{\ov\pi}$
and the multi-set of neighbors of vertex~$i$ in~$\mu(G_{\ov\pi})$.
We refer to the latter difference by the phrase
{\em the contribution of vertex $i$ to the difference
between $G_{\ov\pi}$ and $\mu(G_{\ov\pi})$}.

As a warm-up, we first show that each element of $T$
contributes a non-zero number of units to the difference
(between $G_{\ov\pi}$ and $\mu(G_{\ov\pi})$)
with probability $1-O(\poly(d)/n)^{d/6}$ over the choice of $\ov\pi$.
%
%
%
%
%
%
%
%
Recalling that $i\in T$ contributes to the difference
(between $G_{\ov\pi}$ and $\mu(G_{\ov\pi})$)
if the multi-sets of its neighbors in $G_{\ov\pi}$
and $\mu(G_{\ov\pi})$ differ,
it follows that $i\in T$ contributes to the difference if
{\em for every permutation $\sigma:[2d]\to[2d]$
there exists $j\in[2d]$
such that $g_j(i)\neq\mu(g_{\sigma(j)}(\imu(i)))$}.
%
%
%
%
Hence, the complementary event holds
(i.e., $i$ does not contribute to the difference)
if and only if there exists permutation $\sigma:[2d]\to[2d]$
such that for every $j\in[2d]$
it holds that $g_j(i)=\mu(g_{\sigma(j)}(\imu(i)))$.
%
%
%
%
%
%
%
Thus, the probability that $i$ does not contribute to the difference
is given by
\begin{eqnarray}
\lefteqn{\prob_{\ov\pi}
  \left[\exists\sigma\!\in\!\Sym_{2d}\;
        (\forall j\!\in\![2d])\;\;g_j(i)=\mu(g_{\sigma(j)}(\imu(i)))\right]}
\nonumber \\
&\leq& (2d)!\cdot\max_{\sigma\in\Sym_{2d}}
  \left\{\prob_{\ov\pi}\left[(\forall j\!\in\![2d])\;\;
            g_j(i)=\mu(g_{\sigma(j)}(\imu(i)))\right] \right\}.
\label{true-probability:eq}
\end{eqnarray}
Fixing $\sigma$ that maximizes the probability,
and denoting it $\sigma_i$,
%
%
consider any $J_i\subseteq[d]$ such that for the $j$'s in $J_i$
the multi-sets $\{j,\ceil{\sigma_i(2j)/2}\}$'s are disjoint
(i.e., $\{j,\ceil{\sigma_i(2j)/2}\}\cap\{k,\ceil{\sigma_i(2k)/2}\}=\emptyset$
for any $j\neq k\in J_i$).
Note that we may select $J_i$ such that $|J_i|\geq d/6$,
since taking $j$ to $J_i$ rules out taking (to $J_i$)
at most five other values (i.e., letting $v\eqdef\ceil{\sigma_i(2j)/2}$,
the value $k$ is ruled out if either $k\in\{j,v\}$
or $\sigma_i(2k)\in\{2j-1,2j,2v-1,2v\}$).
%
%
%
%
%
%
%
%















Using this feature of $J_i$, we prove~--

\setcounter{subsection}{1}

\Bcm{\em(Warm-up):}\footnote{One may obtain a better bound
of $O(d/n)^{2d}$ by analyzing \eqref{true-probability:eq} directly,
by considering all the $2d$ events and accounting for their small dependency.
On the other hand, we can obtain higher robustness parameter
by considering smaller sets $J_i$'s (say of size $d/12$),
which suffice for counting vertices that
contribute (say) $d/2$ units to the difference
between $G_{\ov\pi}$ and $\mu(G_{\ov\pi})$.\label{random-bd:fn}}
%
%
%
\label{random-works-warm-up:clm}
\eqref{true-probability:eq} is upper-bounded
by $(2d)^{2d}\cdot(3/n)^{|J_i|}$.
\Ecm

\Bpf
Recalling that $\sigma_i$ maximizes \eqref{true-probability:eq},
we upper-bound \eqref{true-probability:eq} by

\begin{eqnarray}
\nonumber
\lefteqn{(2d)!\cdot
  \prob_{\ov\pi}\left[(\forall j\!\in\! J_i)\;\;
	    g_{2j}(i)=\mu(g_{\sigma_i(2j)}(\imu(i)))\right]} \\
&=& (2d)!\cdot \prod_{j\in J_i}
     \prob_{\pi_{j},\pi_{\ceil{\sigma_i(2j)/2}}}
          \left[g_{2j}(i)=\mu(g_{\sigma_i(2j)}(\imu(i))) \right]
\label{crude-bound:eq}
\end{eqnarray}
where the equality uses the disjointness of
the multi-sets $\{j,\ceil{\sigma_i(2j)/2}\}$
for the $j$'s in $J_i$.
Next, letting $\sigma'_{i}(2j)\eqdef{\ceil{\sigma_{i}(2j)/2}}$,
and $\sigma''_{i}(2j)\eqdef(-1)^{\sigma_{i}(2j)\bmod2}$
(so that $\sigma_{i}(2j)=2\sigma'_{i}(2j)-0.5(1-\sigma''_{i}(2j))$),
we upper-bound \eqref{crude-bound:eq} by
%
\begin{equation} \label{crude-bound:eq2}
(2d)!\cdot \prod_{j\in J_i}
	\prob_{\pi_{j},\pi_{\sigma'_i(2j)}}\left[
	\pi_{j}(i)=\mu(\pi^{\sigma''_i(2j)}_{\sigma'_i(2j)}(\imu(i)))
           \right]
\;<\; (2d)^{2d}\cdot(3/n)^{|J_i|},
\end{equation}
%
%
%
where $\prob_{\pi_j,\pi_j}[\cdot]$ stands for $\prob_{\pi_j}[\cdot]$
and $\pi^1$ stands for $\pi$,
while the inequality is justified by considering
the following three cases (w.r.t each $j\in J_i$).
\BE
%
%
\item If $k\eqdef\sigma'_i(2j)\neq j$
(equiv., $\sigma_i(2j)\not\in\{2j-1,2j\}$),
then, letting $b=\sigma''_i(2j)$,
the corresponding factor in the l.h.s of \eqref{crude-bound:eq2} is
$$\prob_{\pi_{j},\pi_{k}}\left[\pi_{j}(i) = \mu(\pi^b_k(\imu(i)))\right]$$
which equals $1/n$ by fixing $\pi_k$, letting $v=\mu(\pi^b_k(\imu(i)))$,
and using $\prob_{\pi_j}[\pi_j(i)\!=\!v]=1/n$.
\item If $\sigma_i(2j)=2j$,
then the corresponding factor in the l.h.s of \eqref{crude-bound:eq2} is
$$\prob_{\pi_{j}}\left[\pi_{j}(i) = \mu(\pi_j(\imu(i)))\right]$$
which is at most $1/(n-1)$ since $\mu(i)\neq i$;
specifically, fixing the value of $\pi_j(\imu(i))$,
and denoting this value by $v$,
%
%
leaves $\pi_j(i)$ uniformly distributed in $[n]\setminus\{v\}$,
which means that
$\prob_{\pi_j}[\pi_j(i)\!=\!\mu(v)\,|\,v\!=\!\pi_j(\imu(i))]\leq1/(n-1)$
(where equality holds if $\mu(v)\neq v$).
%
%
%
%
\item If $\sigma_i(2j)=2j-1$,
then the corresponding factor in the l.h.s of \eqref{crude-bound:eq2} is
$$\prob_{\pi_{j}}\left[\pi_{j}(i) = \mu(\pi_j^{-1}(\imu(i)))\right]$$
which is shown to be less than $3/n$.
%
%
In this case, we consider two sub-cases
depending on whether or not $\pi_j(i) = \imu(i)$,
while noting that the first case occurs with probability $1/n$ whereas
$\prob_{\pi_j}[\pi_{j}(i)=\mu(\pi_j^{-1}(\imu(i)))|\pi_j(i)\neq\imu(i)]
 \leq1/(n-1)$.
%
%
%
%
%
%
%
%
%
%
%
%
%
%
\EE

Hence, each of the factors in the l.h.s of \eqref{crude-bound:eq2}
is upper-bounded by $3/n$, and the claim follows.
\Epf

\mypar{The general case.}

%
The same argument generalizes to any set $I\subseteq T$
such that $I\cap\mu(I)=\emptyset$.
In such a case we get





















\begin{eqnarray}
\lefteqn{\prob_{\ov\pi}\left[(\forall i\!\in\!I)\,
                    (\exists\sigma_i\!\in\!\Sym_{2d})\,
                    (\forall j\!\in\![2d])\;\;
                  g_j(i)=\mu(g_{\sigma_i(j)}(\imu(i)))\right]}
\nonumber \\
&\leq& (2d)!^{|I|}\cdot\max_{\sigma_1,\ldots,\sigma_n}
  \left\{\prob_{\ov\pi}\left[(\forall i\!\in\!I)\,(\forall j\!\in\![2d])
            \;\;g_j(i)=\mu(g_{\sigma_i(j)}(\imu(i)))\right] \right\}
\label{true-probability:eq2}
\end{eqnarray}
\Bcm{\em(Actual analysis):}
\label{random-works-actual:clm}
\eqref{true-probability:eq2} is upper-bounded by
%
\begin{equation} \label{true-probability:eq3}
(2d)^{2d\cdot|I|}\cdot(3/(n-2(|I|-1)))^{|I|\cdot d/6}.
\end{equation}
\Ecm

\Bpf
Fixing $\sigma_1,....,\sigma_n$
that maximize \eqref{true-probability:eq2},
we proceed as in the proof of Claim~\ref{random-works-warm-up:clm}.
For every $i\in I=\{i_1,\ldots,i_m\}$,

we fix a set $J_i$ of size at least $d/6$ such that
the multi-sets $\{j,\ceil{\sigma_i(2j)/2}\}$'s are disjoint,
and refer to the events $E_{j,k}(\pi_1,\ldots,\pi_{2d})$
that depend only on the value of $\pi_j$
and $\pi^{\sigma''_{i_k}\!(2j)}_{\sigma'_{i_k}\!(2j)}$
on the points $i_1,\ldots,i_{k-1}$ and $\imu(i_1),\ldots,\imu(i_{k-1})$,
respectively, where $\sigma'_{i}(2j)\eqdef{\ceil{\sigma_{i}(2j)/2}}$,
and $\sigma''_{i}(2j)\eqdef(-1)^{\sigma_{i}(2j)\bmod2}$
(as in the proof of Claim~\ref{random-works-warm-up:clm}).
Specifically, $E_{j,k}(\pi_1,\ldots,\pi_{2d})$ is the event
$$(\forall k'\!\in\![k-1])
             \;\;g_{2j}(i_{k'})=\mu(g_{\sigma_{i_{k'}}\!(2j)}(\imu(i_{k'})))$$
which can be written as
$$(\forall k'\!\in\![k-1]) \;\; \pi_j(i_{k'})
  =\mu(\pi^{\sigma''_{i_{k'}}\!(2j)}_{\sigma'_{i_{k'}}\!(2j)}(\imu(i_{k'}))).$$
Using the disjointness of the $\{j,\ceil{\sigma_i(2j)/2}\}$'s in $J_i$,
we upper-bound \eqref{true-probability:eq2} by

%
%
%
\begin{eqnarray}
\lefteqn{(2d)!^{m}\cdot\prod_{k\in[m]~} \prod_{~j\in J_{i_k}}
     \prob_{\pi_1,\ldots,\pi_{2d}}
          \left[g_{2j}(i_k)=\mu(g_{\sigma_{i_k}\!(2j)}(\imu(i_k)))\left|
                E_{j,k}(\pi_1,....,\pi_{2d}) \right.\right]}
\nonumber \\
&=& (2d)!^{m}\cdot\prod_{k\in[m]~} \prod_{~j\in J_{i_k}}
     \prob_{\pi_1,\ldots,\pi_{2d}}
          \left[\pi_{j}(i_k)
      =\mu(\pi^{\sigma''_{i_k}\!(2j)}_{\sigma'_{i_k}\!(2j)}
	  (\imu(i_k)))\left|E_{j,k}(\pi_1,....,\pi_{2d}) \right.\right] .
\label{true-probability:eq4}
%
%
\end{eqnarray}
%
%
%
%
%
%
%
%
%
%
%
%
%
%
Now, when analyzing the foregoing conditional probability
in \eqref{true-probability:eq4},
we consider two cases for each $k\in[m]$ and $j\in J_{i_k}$.
If $j\neq\sigma'_{i_k}\!(2j)$, then
%
we fix the value of each of these two permutations
(i.e., $\pi_j$ and $\pi_{\sigma'_{i_k}\!(2j)}$)
on the corresponding $k-1$ points that occur in the condition $E_{j,k}$,
and the value of these permutations on the $k^\xth$ points
(i.e., $i_k$ and $\imu(i_k)$) is restricted accordingly
(i.e., to the remaining $n-(k-1)$ values).
%
Otherwise (i.e., $j=\sigma'_{i_k}\!(2j)$),
we fix the value of $\pi_j$ on these $2(k-1)$ points.
Hence, the argument in the warm-up analysis applies with $n$
replaces by either $n-(k-1)$ or $n-2(k-1)$.
It follows that \eqref{true-probability:eq4} is upper-bounded by
$$(2d)!^{m}\cdot\prod_{k\in[m]~}(3/(n-2(m-1)))^{|J_{i_k}|}.$$
Using $|J_{i_k}|\geq d/6$ for every $k\in[m]$,

the claim follows.
\Epf



\medskip

Recall that \eqref{true-probability:eq3} refers to
a fixed set $I\subseteq T$ such that $I\cap\mu(I)=\emptyset$,
and that it constitutes an upper bound
on the probability (over the choice of $\ov\pi$)
that, for each $i\in I$ there exists
a permutation $\sigma_i:[2d]\to[2d]$
such that $g_j(i)=\mu(g_{\sigma_i(j)}(\imu(i)))$
holds for all $j\in[2d]$.
This upper bound
(i.e., $(2d)^{2d\cdot|I|}\cdot(3/(n-2(|I|-1)))^{|I|\cdot d/6}$)
simplifies to $(2d)^{2d\cdot|I|}\cdot(5/n)^{|I|\cdot d/6}$,
provided that $|I|\leq\ceil{n/6}$.

%
%
%
%
%
%
%

Recalling that $t\eqdef|T|\in[n]$,
we shall upper-bound the probability (over the choice of $\ov\pi$)
that $T$ contains a $\ceil{t/2}$-subset $T'$
such that for each $i\in T'$ there exists
a permutation $\sigma_i:[2d]\to[2d]$
such that $g_j(i)=\mu(g_{\sigma_i(j)}(\imu(i)))$
holds for all $j\in[2d]$.
We do so by taking a union bound over all $\ceil{t/6}$-subsets $I$
such that $I\cap\mu(I)=\emptyset$
and for each $i\in I$ there exists a permutation $\sigma_i:[2d]\to[2d]$
such that $g_j(i)=\mu(g_{\sigma_i(j)}(\imu(i)))$ holds for all $j\in[2d]$.
(Note that such a $\ceil{t/6}$-subset $I$ exists
in each $\ceil{t/2}$-subset $T'$, and that $\ceil{t/6}<n/3$.)
%
%
Using the aforementioned simplified form of \eqref{true-probability:eq3},
we conclude that, with probability at most
%
$$\binom{t}{\ceil{t/6}}
   \cdot(2d)^{2d\cdot\ceil{t/6}}\cdot(5/n)^{\ceil{t/6}\cdot{d/6}}
   \;<\; 2^t\cdot(5\cdot(2d)^{12}/n)^{\ceil{t/6}\cdot d/6}
   \;=\; \exp(-\Omega(dt\log n))$$
over the choice of $\ov\pi$,
the set $T$ contains no $\ceil{t/6}$-subset $I$ as above.
This means that, with probability at most $\exp(-\Omega(d t\log n))$,
less than ${t/2}$ of the indices $i\in T$ contribute
a non-zero number of units to the difference
(between $G_{\ov \pi}$ and $\mu(G_{\ov \pi})$).


Letting $c'=1/2$ and
considering all (non-trivial) permutations $\mu:[n]\to[n]$,
we conclude that the probability, over the choice of $\ov\pi$,
that $G_{\ov\pi}$ is not $c'$-robustly self-ordered is at most
\begin{eqnarray*}
\sum_{t\in[n]}\binom{n}{t}\cdot\exp(-\Omega(d t\log n))
&=& \sum_{t\in[n]}\exp(-\Omega((d-O(1))\cdot t\log n)) \\
&=& \exp(-\Omega((d-O(1))\cdot \log n)),
\end{eqnarray*}
and the claim follows for the permutation model
(and for any sufficiently large $d$).
%

As stated upfront,
using the general result of~\cite[Thm.~1.3]{GJKW},
%
%
%
%
we infer that a uniformly distributed $2d$-regular $n$-vertex
multi-graph fails to be $c'$-robustly self-ordered with probability $o(1)$.
Lastly, recalling that such a $2d$-regular multi-graph is actually
a simple graph with probability $\exp(-((2d)^2-1)/4)$,
the theorem follows.
%
\EPF

\paragraph{Digest.}
The proof of Theorem~\ref{random-works:thm} is quite similar
to the proof Claim~4.1.1, but it faces
two complications that were avoided in the prior proof
(by using edge-colors and implicitly directed edges).
Most importantly, the current proof has to handle equality
between multi-sets instead of equality between sequences.
This is done by considering all possible ordering of these multi-sets,
which amounts to taking a union bound over all possible ordering
and results in more complicated analysis and notation.
(Specifically, see the introduction of $\sigma_i$'s and $J_i$'s
and the three cases analyzed in the warm-up.)
In addition, since edges are defined by permutations over the vertex-set
rather than by perfect matching, we have to consider both the forward
and backward direction of each permutation,
which results in further complicating the analysis and the notation.
(Specifically, see the introduction of $\sigma'_i$'s and $\sigma''_i$'s
and the three cases analyzed in the warm-up.)










































































































































































































































































































































































































































































































































































































































































































































































































































































































































































































































































































































































































































































































































































































































































































































































































































































































































































































































































































































































































































































































































































































































































































































































































































































































































































































































































































































































































































































































































































































































































































































































































































































































































